\documentclass[a4paper,12pt]{report}
\usepackage[competences,niveau={1MA1},nfiche={}, annee={Emilie-Gourd, 2024--2025}, auteur={},theme={Exercices aléatoirisés -- ER juin}]{packages/bandeaux}
\usepackage{packages/boites}
\usepackage{packages/fontes}
\usepackage{packages/courslatex}
\setlist[itemize]{align=left,labelsep=1em,leftmargin=*,itemsep=5pt,parsep=5pt,topsep=5pt,rightmargin=0cm}



\begin{document}

Pour commencer les révisions~:

\begin{center}
	\qrwithlabel{Notions de base sur les angles}{https://coopmaths.fr/alea/?uuid=dc8c9&id=9ES2-9&uuid=19812&id=11ES2-2&uuid=d12db&id=11ES2-1&lang=fr-CH&v=eleve&es=1011001&title=Angles+-+notions+de+base}
	\qrwithlabel{Relations métriques dans un triangle rectangle}{https://coopmaths.fr/alea/?uuid=4f81e&id=1mG3-6&n=10&d=10&s=13&s2=true&s3=true&cd=0&alea=ozGW&lang=fr-CH&v=eleve&title=Relation+metriques&es=0211001}
	\qrwithlabel{Formules aire et périmètre disque}{https://coopmaths.fr/alea/?uuid=ff386&id=10GM1-3&lang=fr-CH&v=eleve&es=2011001&title=Fomule+disque}

	\medskip
\qrwithlabel{Trigonométrie de base}{https://coopmaths.fr/alea/?uuid=c6b9c&id=1mT-2&uuid=0ac11&id=1mT-3&n=2&d=10&s=true&cd=0&uuid=0ac11&id=1mT-3&n=2&d=10&s=false&cd=0&uuid=bd6b1&id=1mT-5&n=3&d=10&s=true&s2=7&s3=1&cd=0&uuid=bd6b1&id=1mT-5&n=3&d=10&s=false&s2=7&s3=1&cd=0&uuid=35e0b&id=1mT-6&n=2&d=10&s=false&cd=1&lang=fr-CH&v=eleve&es=1011001&title=Trigo+de+base}
	\qrwithlabel{Problèmes de trigonométrie}{https://coopmaths.fr/alea/?uuid=2045e&id=1mT-9&n=5&d=10&s=false&s2=1-3-4-5-6&cd=1&lang=fr-CH&v=eleve&es=3211001&title=Problemes+de+trigo}
\end{center}

\end{document}
